% Introdução (versão compacta para o limite de 10 páginas). Versão completa: artigo/completo/introducao.tex.
% As hipóteses foram escritas ANTES de qualquer análise (Issue #15, Lab03S01); aqui só foram condensadas, sem
% mudar a expectativa de cada uma. A redação original está em artigo/completo/introducao.tex.
\section{Introdução}

As métricas DORA (\textit{DevOps Research and Assessment}) tornaram-se o padrão de mercado para medir o
desempenho de entrega de software~\cite{forsgren2018accelerate}. Quatro medidas descrevem duas dimensões:
velocidade (\textit{deployment frequency} e \textit{lead time for changes}) e estabilidade (\textit{change failure
rate} e tempo de recuperação). O relatório \textit{Accelerate State of DevOps} sustenta que as melhores equipes
vão bem nas duas dimensões, isto é, que velocidade e estabilidade não são um \textit{trade-off}.

O GitHub não registra ``deploys em produção'' nem ``falhas em produção''; só é possível \emph{inferir} as métricas
por \textbf{proxies}: uma \textit{release} publicada faz o papel de deploy, e uma execução de CI que falhou, o de
falha. Por isso cada métrica minerada precisa de uma medida de confiança. Este trabalho \textbf{calcula as métricas DORA} de repositórios populares que usam GitHub Actions e
\textbf{mede o quanto as conclusões dependem das definições operacionais}. Hipóteses informais, escritas antes
de ver os dados:

\begin{description}\setlength{\itemsep}{1pt}
  \item[H1 (RQ01).] Distribuição muito assimétrica, com mediana entre uma \textit{release} a cada duas semanas e
  uma por mês (\textit{Medium}) e cauda de projetos com mais de uma por semana.
  \item[H2 (RQ02).] A mediana por commit (b) será bem menor que a por \textit{release} (a), pois um commit antigo
  incluído na \textit{release} faz (a) crescer e afeta pouco (b); \textit{lead times} medianos de semanas a poucos meses.
  \item[H3 (RQ03).] O CFR (a), de CI, será maior que o (b), de \textit{releases} corretivas, com mediana de (a) entre
  10\% e 30\%: falhas de pipeline são comuns e raramente chegam a uma \textit{release} corretiva.
  \item[H4 (RQ04).] A maioria dos episódios se resolve em horas, com cauda longa; mediana por repositório entre 1 hora
  e 1 dia, e parcela relevante de episódios censurados.
  \item[H5 (RQ05).] Correlação de Spearman fraca ou levemente negativa entre frequência e CFR, em linha com a tese
  DORA, mas com ruído por causa dos proxies.
  \item[H6 (RQ06).] Linguagem e tipo de projeto terão associação mais forte que a popularidade; aplicações publicam
  mais que bibliotecas; idade e contribuidores terão efeitos pequenos, e a maioria das diferenças perderá
  significância após a correção para comparações múltiplas.
  \item[H7 (RQ07).] Mais de 30\% dos repositórios mudarão de categoria ao trocar a definição de deploy e as variantes
  de LT e CFR, com $\kappa$ ponderado no máximo moderado.
\end{description}
